% ====================================================================
% TEMA 4 — JOSÉ ÁNGEL
% Innovación y Propuestas de Reclutamiento
% ====================================================================

\section{Innovación y Propuestas de Reclutamiento}

\subsection{Nuevas tendencias de reclutamiento: modelos 2.0 y 3.0}

La digitalización ha transformado profundamente el reclutamiento. La teoría distingue el
\textbf{e-recruitment} o \textbf{reclutamiento 2.0}, basado en el uso de internet y portales
de empleo, del \textbf{reclutamiento 3.0} o \textit{social recruiting}, definido como «el
proceso por el cual las empresas buscan e identifican los posibles candidatos para cubrir
sus necesidades internas de empleo a través de las redes sociales multiplataformas de
internet» [8].

CaixaBank combina ambos modelos de forma madura:

\begin{itemize}
  \item \textbf{Reclutamiento 2.0.} La entidad dispone de un portal propio,
        \textit{CaixaBankCareers}, donde los candidatos registran su candidatura de forma
        digital, acceden a información sobre el proceso y realizan las pruebas en línea
        [1]. Además, colabora con agencias de trabajo temporal como
        Randstad para ampliar el alcance externo del proceso.

  \item \textbf{Reclutamiento 3.0.} CaixaBank mantiene una presencia activa en LinkedIn
        —con más de 195\,000 seguidores— donde publica ofertas, difunde contenido de
        \textit{employer branding} y contacta directamente con perfiles pasivos
        [5]. En 2019 lanzó el \textbf{People Xperience Hub},
        una iniciativa de \textit{Recruitment Process Outsourcing} (RPO) que estandariza
        y digitaliza todo el ciclo de atracción de talento [6].
        También utiliza Instagram y otras redes sociales para reforzar su marca de
        empleador.
\end{itemize}

Esta estrategia presenta las ventajas señaladas en la teoría: menor coste respecto al
reclutamiento tradicional, mayor alcance geográfico y mejor capacidad de precalificación
y filtrado de candidatos. No obstante, CaixaBank debe gestionar con cuidado el riesgo de
reputación \textit{online} y el cumplimiento normativo en el tratamiento de datos
personales de los candidatos.

\subsection{Uso e impacto del Vídeo-CV}

El \textbf{videocurrículum} es una tendencia consolidada en los procesos modernos de
reclutamiento: el candidato se graba presentando su perfil, permitiendo al equipo de
selección valorar competencias comunicativas y actitudinales que escapan al CV
tradicional [8].

El proceso de selección de CaixaBank consta de seis fases digitalizadas [1], lo que
lo hace compatible con la incorporación del videocurrículum como herramienta de
preselección. Su estructura online —con pruebas y dinámicas realizadas en remoto—
encaja directamente con el modelo que describe la teoría.

El impacto de esta herramienta sería doble. Para la empresa, reduce desplazamientos y
costes logísticos, acelera la criba y permite revisar las grabaciones de forma asíncrona.
Para el candidato, ofrece mayor flexibilidad horaria y elimina barreras geográficas,
aunque puede introducir sesgos si no se establecen criterios de evaluación estructurados.
Un riesgo adicional es la brecha digital: candidatos sin acceso a tecnología adecuada
pueden quedar en desventaja, lo que puede comprometer la equidad del proceso.

\subsection{Propuestas de mejora para el reclutamiento futuro}

\subsubsection{Mejoras a corto plazo}

\begin{itemize}
  \item \textbf{Incorporación de Inteligencia Artificial en el cribado.} La teoría
        recoge que herramientas de \textit{machine learning} permiten filtrar currículos
        automáticamente y realizar búsquedas semánticas para detectar candidatos idóneos
        [8]. CaixaBank podría integrar un sistema ATS (\textit{Applicant
        Tracking System}) con IA para reducir el tiempo de criba y minimizar sesgos
        inconscientes, garantizando siempre la supervisión humana de las decisiones
        finales.

  \item \textbf{Gamificación del proceso de selección.} La introducción de pruebas
        situacionales en formato de juego (\textit{serious games}) durante la fase de
        preselección permitiría evaluar competencias como la toma de decisiones o el
        trabajo en equipo de forma más objetiva y atractiva para los candidatos,
        mejorando a la vez la experiencia de candidatura (\textit{candidate experience}).
\end{itemize}

\subsubsection{Mejoras a medio plazo}

\begin{itemize}
  \item \textbf{Potenciación del \textit{employer branding} en redes sociales.}
        Aunque CaixaBank ya cuenta con presencia en LinkedIn, podría intensificar
        el uso de contenido generado por empleados (\textit{employee advocacy}) y
        la difusión de testimonios reales para atraer talento pasivo, reduciendo
        la dependencia de canales de pago y mejorando su posicionamiento como
        empleador de referencia en el sector financiero.

  \item \textbf{Programas de reclutamiento universitario (\textit{campus recruiting}).}
        La teoría señala a las instituciones docentes como fuente clave para obtener
        candidatos cualificados [8]. CaixaBank podría estructurar
        acuerdos estables con universidades —más allá de las prácticas actuales— para
        identificar talento en etapas tempranas y construir una bolsa de candidatos
        preseleccionados para posiciones de perfil junior.
\end{itemize}

\subsubsection{Indicadores de seguimiento propuestos}

Para evaluar la eficacia de las mejoras, se propone monitorizar los siguientes
indicadores, coherentes con los ratios de eficiencia del reclutamiento definidos
en la teoría [8]:

\begin{itemize}
  \item \textbf{\textit{Yield ratio}}: proporción de candidatos que superan cada
        fase del proceso, para identificar cuellos de botella.
  \item \textbf{\textit{Time-to-hire}}: tiempo medio desde la publicación de la
        oferta hasta la contratación, con el objetivo de reducirlo por debajo de
        los 41 días actuales.
  \item \textbf{Tasa de aceptación de oferta}: porcentaje de candidatos que aceptan
        la propuesta, indicador de la percepción de la marca empleadora.
  \item \textbf{Satisfacción del candidato}: encuesta post-proceso para medir la
        \textit{candidate experience} e identificar áreas de mejora.
\end{itemize}
